\chapter{Structural Regions}
\label{ch:regions}

Chapter~\ref{ch:ontology} introduced the text object as a distinct kind of entity from a motion: where a motion describes a path the cursor could travel, a text object describes a region directly, defined by the structure it belongs to rather than by any particular route to its edges. Chapter~\ref{ch:language} placed text objects in the grammar, showing that they fill the same slot a motion fills when paired with an operator. This chapter returns to text objects on their own terms, cataloging the structural regions Vim recognizes and examining the inner/around distinction that runs through nearly all of them.

\section{The Inner/Around Distinction}

Almost every text object Vim defines comes in two forms, distinguished by a prefix letter: \texttt{i} for "inner," selecting a structure's content without its delimiters, and \texttt{a} for "around," selecting the structure together with its delimiters and, in most cases, one adjacent span of trailing whitespace. \texttt{ci(} changes the contents of a pair of parentheses without disturbing the parentheses themselves; \texttt{ca(} changes the parentheses and their contents together, removing them entirely if nothing is typed to replace them. This single distinction, learned once, applies almost without exception across every structural region below, which is precisely the kind of multiplicative economy Section~\ref{sec:multiplication} in the previous chapter described: understanding \texttt{i}/\texttt{a} as a general rule, rather than as forty separate memorized command pairs, is what makes the whole catalog tractable.

\section{Word Regions}

\texttt{iw} selects the inner word under or after the cursor: the run of alphanumeric or punctuation characters the cursor sits within, using the same word-boundary rules as the \texttt{w} motion from Chapter~\ref{ch:coordinates}. \texttt{aw} selects the word together with one adjacent span of whitespace, generally the trailing whitespace if any exists, or the leading whitespace otherwise. \texttt{iW} and \texttt{aW} apply the same inner/around distinction to Vim's coarser, punctuation-insensitive WORD unit. The practical effect of the whitespace-inclusion rule in the around forms is that repeatedly deleting a word with \texttt{daw} does not slowly accumulate stray double spaces the way repeatedly deleting with \texttt{diw} would; the around form is generally the correct default for actually removing a word from a sentence, while the inner form is generally correct when the word is about to be replaced with something of a different length via \texttt{ciw}.

\section{Quote Regions}

\texttt{i"}, \texttt{i'}, and \texttt{i\`{}} select the contents of the nearest double-quoted, single-quoted, or backtick-quoted string on the current line, without the quote marks; \texttt{a"}, \texttt{a'}, and \texttt{a\`{}} include the quote marks and, where present, one trailing space. Quote text objects are line-bound in classic Vim: they search only within the current line for the nearest enclosing or following quoted pair, which is a deliberate limitation rather than an oversight, since quote characters are common enough that searching an entire buffer for a matching pair could easily select a much larger, unintended span. A cursor positioned before the first quote character on a line will target the next quoted string that appears later on that same line, which is worth knowing since it means these text objects do not strictly require the cursor to already be inside the quotes to function.

\section{Bracket and Parenthesis Regions}

\texttt{i(} and \texttt{i)} (equivalently, \texttt{ib}, provided as a shorter mnemonic) select the contents of the nearest enclosing pair of parentheses; \texttt{i\{}, \texttt{i\}}, and the mnemonic \texttt{iB} do the same for curly braces; \texttt{i[} and \texttt{i]} for square brackets; \texttt{i<} and \texttt{i>} for angle brackets. The corresponding \texttt{a} forms include the delimiters themselves. Unlike quote regions, bracket regions are not restricted to the current line and correctly handle nesting: invoked with the cursor inside a nested pair of parentheses, \texttt{i(} selects the innermost enclosing pair, and pressing the same command again while that selection is active (in Visual mode, introduced in Chapter~\ref{ch:ontology} and taken up further throughout Part III) expands the selection outward to the next enclosing pair, a behavior that makes working with deeply nested expressions, common in Lisp-family languages and in deeply nested JSON or configuration data, considerably more tractable than manually counting matched delimiters by eye.

\section{Tag Regions}

\texttt{it} selects the content between an opening and closing markup tag, such as the text between \texttt{<p>} and \texttt{</p>} in HTML or between an opening and closing element in XML, without the tags themselves; \texttt{at} includes the tags. This text object is markup-aware in a way the bracket text objects are not: it parses tag names and matches an opening tag to its corresponding closing tag by name, correctly skipping over unrelated intervening tags at different nesting depths, rather than relying on any single, universal delimiter character the way parenthesis matching does. It is consequently one of the more computationally involved text objects Vim provides, and one of the more immediately useful ones for anyone editing HTML, XML, or similar markup on a regular basis, covered further in the context of writing workflows in Chapter~20.

\section{Sentence and Paragraph Regions}

\texttt{is} and \texttt{as} select the current sentence, using the same sentence-boundary rules as the \texttt{(} and \texttt{)} motions from Chapter~\ref{ch:coordinates}, with the around form including trailing whitespace; \texttt{ip} and \texttt{ap} do the same for the current paragraph, using the same blank-line-delimited definition the \texttt{\{} and \texttt{\}} motions use. These text objects are the direct region-based counterparts of the sentence and paragraph motions already covered, and the relationship between a coordinate system in Chapter~\ref{ch:coordinates} and the corresponding structural region here is not a coincidence: wherever a motion moves to the boundary of some structural unit, that unit is generally also available as a text object addressing its interior directly, which is a pattern worth watching for throughout the rest of the book rather than treating each pairing as a separate fact to memorize.

\section{Block Regions}

Visual Block mode, taken up further in Part III, is itself a way of selecting a rectangular region of a buffer directly, by column range across multiple lines, rather than a text object in the operator-plus-target sense used elsewhere in this chapter. It is mentioned here because it fills a structurally similar role: a way of designating a region other than a simple linear span of characters. Where the text objects above select spans defined by delimiters, word boundaries, or blank lines, block selection selects a span defined by column position, useful for editing aligned columns of data or adding the same prefix to several consecutive lines at once.

\section{User-Defined Regions}

Everything covered so far in this chapter is built into Vim, but the inner/around pattern is not closed to further extension. Vimscript (Part VI) permits defining entirely new text objects through custom operator-pending mappings, letting a user or a plugin author extend the same \texttt{i}/\texttt{a} grammar to structures Vim has no built-in notion of, such as an entire function definition in a specific programming language, a Markdown heading section, or an indentation level. This book, consistent with the plugin-free philosophy stated in the Preface, does not cover any specific third-party text-object definitions, but it is worth knowing the mechanism exists, because it is direct evidence for a claim made repeatedly throughout this chapter: the inner/around distinction is not a fixed, closed list of special cases Vim happens to support, but a general grammatical pattern, available to be extended to any structural unit a user can define programmatically, using exactly the same two-letter vocabulary already learned for the built-in cases above.

\section{Structural Regions and the Grammar}

It is worth closing this chapter, and with it Part II, by returning explicitly to where structural regions sit within the larger rule established in Chapter~\ref{ch:language}. A text object is never used alone in ordinary editing; it is always the target half of an operator-target pair, or the boundary of a Visual-mode selection about to be acted on. Everything catalogued in this chapter multiplies against everything catalogued in the next chapter's treatment of operators, exactly as Section~\ref{sec:multiplication} of Chapter~\ref{ch:language} predicted it would: roughly a dozen structural regions, each with an inner and around form, multiplied against a handful of core operators, already yields several dozen distinct, meaningful commands, the overwhelming majority of which this book will never explicitly demonstrate, because by this point in the book they no longer need to be. That is the payoff Part II was written to deliver.
